\documentclass[12pt]{article}
\usepackage{amsmath}
\usepackage{amssymb}
\usepackage{amsthm}
\usepackage{geometry}
\geometry{margin=1.4in}
\usepackage{hyperref}
\usepackage{setspace}
\setstretch{1.4}

\newtheorem{theorem}{Theorem}
\newtheorem{conjecture}[theorem]{Conjecture}
\theoremstyle{definition}
\newtheorem{definition}[theorem]{Definition}
\theoremstyle{remark}
\newtheorem{remark}[theorem]{Remark}

\title{The Resonance Object}
\author{Diego Rinc\'on\\
\small Independent}
\date{}

\begin{document}

\maketitle

\begin{abstract}
A wall is not transparent everywhere.
It is transparent at a frequency.
The resonance object is the structure that finds the frequency.
For the Riemann Hypothesis, the resonance object has not been found.
This paper is about what it would have to be.
\end{abstract}

\section{What Resonance Means}

A wall is not an end.
It is a mismatch.

The current grammar approaches the problem at a certain frequency ---
a certain register of mathematical structure.
The wall is where that frequency stops passing through.
Not where mathematics stops.
Where this mathematics stops.

The resonance method: find the frequency the wall is transparent at.
Change the register. Approach from the right angle.

This is not a vague prescription.
It has a mathematical content:
the wall is a grammatical boundary \cite{grammar},
a place where the current cohomological grammar
cannot construct the required coboundary.
The resonance object is a structure in a new or extended grammar
that makes the coboundary constructible.

The wall does not move.
The grammar expands to include what the wall was made of.

\section{Historical Resonances}

The method has worked before.

\textbf{Wiles and BSD for rank 0--1.}
The wall for BSD at rank $\leq 1$ was crossed
by the resonance of two grammars:
modular forms (automorphic) and elliptic curves (arithmetic).
The resonance object was the modular parametrization $X_0(N) \to E$.
Once this object existed,
the Euler system could be constructed (Kolyvagin),
and the wall fell for $r \leq 1$.
The resonance object was Shimura--Taniyama--Weil,
proved by Wiles.

\textbf{Deligne and the Weil conjectures.}
The wall was the Riemann Hypothesis over finite fields.
The resonance object was \'etale cohomology:
a new cohomology theory (Grothendieck's grammar)
in which the Frobenius acted with the right eigenvalue properties.
Once the resonance object existed (the \'etale cohomology groups $H^i_{\rm\acute{e}t}$),
the Weil conjectures followed.
The resonance object was built, and then the proof was the coboundary it generated.

\textbf{Perelman and Poincar\'e.}
The wall was the classification of 3-manifolds.
The resonance object was Ricci flow with surgery:
a geometric evolution equation that deformed
any simply connected 3-manifold toward the 3-sphere.
The resonance was between the topology (simply connected)
and the geometry (positive curvature under Ricci flow).
Once the resonance object was constructed (Hamilton's Ricci flow, Perelman's surgery),
the proof was the flow itself.

In each case: a new object, living in a new or extended grammar,
that matched the frequency of the wall and passed through it.

\section{The Candidates for RH}

For the Riemann Hypothesis,
the resonance object has not been found.
Four candidates have been proposed.

\textbf{Candidate 1: The Hilbert--P\'olya operator.}
A self-adjoint operator $H$ on a Hilbert space
whose eigenvalues are $\{1/2 + it : \zeta(1/2+it) = 0\}$.
If $H$ exists, RH follows (self-adjoint operators have real eigenvalues;
the eigenvalues $1/2 + it$ have real part $1/2$).
The resonance: between spectral theory (self-adjointness)
and number theory (zeros of $\zeta$).
The object: not found.
The space it would act on: not identified.

\textbf{Candidate 2: Connes' spectral triple.}
Connes constructed a spectral triple on the ad\`ele class space
$\mathbf{A}_\mathbf{Q}/\mathbf{Q}^*$ \cite{connes}.
The Dirac operator $D$ of this triple has absorption spectrum
related to the zeros of $\zeta$.
The resonance: between noncommutative geometry (spectral triple)
and arithmetic (adèle class space, zeros).
The candidate is serious. The proof that the eigenvalues of $D$
are exactly the zeros, with the right multiplicities, is incomplete.

\textbf{Candidate 3: The $\mathbb{F}_1$-Frobenius.}
If geometry over $\mathbb{F}_1$ can be made rigorous \cite{f1},
the Frobenius of $\mathrm{Spec}\,\mathbb{Z}$ over $\mathbb{F}_1$
would be the resonance object.
The resonance: between the geometry of $\mathbb{F}_1$-curves
(where RH is a theorem) and the number theory of $\zeta(s)$
(where RH is a conjecture).
The candidate is structural.
The grammar does not yet exist.

\textbf{Candidate 4: The derived Hecke algebra.}
Venkatesh proposed that the cohomology of arithmetic manifolds
carries extra structure from a derived Hecke algebra ---
a non-commutative ring acting on cohomology in degree higher than classical.
The resonance: between the derived structure (new algebra)
and the $L$-functions (classical Hecke operators).
If the derived action can be controlled,
it may unlock BSD for $r \geq 3$ first,
and possibly RH by a related mechanism.
The candidate is recent and active.

\section{What the Resonance Object Must Satisfy}

Whatever the resonance object is,
it must satisfy the following:

\begin{enumerate}
\item \textbf{It lives in an extension of the current grammar.}
The current grammar cannot prove RH \cite{grammar}.
The resonance object is in the extended grammar.
It is not a combinatorial trick within the current one.

\item \textbf{It matches the symmetry of $\zeta$.}
The functional equation $\xi(s) = \xi(1-s)$ gives $\mathrm{Re}(s) = 1/2$
as the axis of symmetry.
The resonance object must respect this symmetry:
it must be self-adjoint (Hilbert--P\'olya),
or Frobenius-symmetric ($\mathbb{F}_1$),
or scale-invariant (Connes).

\item \textbf{It is canonical.}
Not a structure chosen to make the proof work.
A structure that was always there,
that the current grammar was too coarse to see.
The \'etale cohomology groups $H^i_{\rm\acute{e}t}$ were always there.
Grothendieck made them visible.

\item \textbf{It generates the coboundary.}
In the cohomological grammar \cite{grammar},
the proof of RH is a coboundary $\eta$
with $d\eta = \omega_{\rm RH}$.
The resonance object is the form $\eta$.
Finding the resonance object and finding the proof are the same thing.
\end{enumerate}

\section{The Question}

Where does the resonance object live?

Not: what is the proof.
The proof follows from the object.
The hard question is the object.

In what space does it live?
In what grammar is it natural?
What does it look like?
Is it an operator? A cohomology class? A curve? A motive? A flow?

The body of papers has been building toward this question.
It has named the wall \cite{six_walls, one_wall}.
It has named the missing symmetry \cite{frobenius}.
It has named the missing grammar \cite{f1}.
It has not found the object.

The object is not in the body.

But the body knows what the object must satisfy.
And that is the most the map can give.

\begin{thebibliography}{9}

\bibitem{grammar}
Rinc\'on, D. (2026).
Mathematics as Cohomological Grammar. Preprint.

\bibitem{connes}
Rinc\'on, D. (2026).
Noncommutative Geometry and the Riemann Hypothesis. Preprint.

\bibitem{f1}
Rinc\'on, D. (2026).
The Field with One Element. Preprint.

\bibitem{frobenius}
Rinc\'on, D. (2026).
The Missing Frobenius. Preprint.

\bibitem{six_walls}
Rinc\'on, D. (2026).
The Six Walls. Preprint.

\bibitem{one_wall}
Rinc\'on, D. (2026).
One Wall: The Local-to-Global Conjecture. Preprint.

\bibitem{riemann2}
Rinc\'on, D. (2026).
The Riemann Hypothesis: A Synthesis. Preprint.

\end{thebibliography}

\end{document}
